\documentclass[11pt]{amsart}

\usepackage[T1]{fontenc}
\usepackage{lmodern}
\usepackage{microtype}
\usepackage{mathtools,amssymb,amsthm}
\usepackage[hidelinks]{hyperref}

\hypersetup{
  pdftitle={A Third-Derivative Counterexample for Chromatic Polynomials}
}

\newtheorem{theorem}{Theorem}
\newtheorem{lemma}[theorem]{Lemma}
\newtheorem{corollary}[theorem]{Corollary}
\theoremstyle{remark}
\newtheorem*{remark}{Remark}

\newcommand{\Lfun}{\mathcal L}
\newcommand{\ThetaG}{\Theta}

\title[A third-derivative counterexample]
{A Third-Derivative Counterexample for Chromatic Polynomials}

\date{}

\subjclass[2020]{05C31, 05C15}
\keywords{chromatic polynomial, logarithmic derivative, generalized theta graph,
counterexample}

\begin{document}

\begin{abstract}
Let $P(G,x)$ be the chromatic polynomial of an $n$-vertex graph $G$.
For $x<0$, set $\Lfun_G(x)=\log\!\bigl((-1)^nP(G,x)\bigr)$.
Conjecture~2 of Dong et al. asserts that
$\Lfun_G^{(k)}(x)<0$ for every graph $G$, every $k\ge2$, and every $x<0$.
We refute it by the simple planar generalized theta graph
$H=\ThetaG(8,8,9,9)$: an exact calculation gives
\[
 \Lfun_H^{(3)}\!\left(-\tfrac12\right)
 =
 \frac{5255579329573016681520193060821596608}
 {18016243181285229121994454910138997625}>0.
\]
The same witness violates the upper inequality in their stronger Conjecture~3.
\end{abstract}

\maketitle

\section{Introduction}

Let $P(G,x)$ denote the chromatic polynomial of a graph $G$ of order $n$.
Its coefficients alternate in sign, so $(-1)^nP(G,x)>0$ for $x<0$ and
\[
 \Lfun_G(x):=\log\!\bigl((-1)^nP(G,x)\bigr)
\]
is well defined on $(-\infty,0)$. Dong, Ge, Gong, Ning, Ouyang, and Tay
proposed the following conjecture \cite[Conjecture~2]{DongEtAl}:
\begin{equation}\label{eq:C2}
 \Lfun_G^{(k)}(x)<0
 \qquad(k\ge2,\ x<0)
\end{equation}
for every graph $G$. They further conjectured that, whenever $G$ is a
noncomplete graph of order $n$ and $Q$ is a chordal proper spanning subgraph
of $G$,
\begin{equation}\label{eq:C3}
 \Lfun_Q^{(k)}(x)
 <\Lfun_G^{(k)}(x)
 <\Lfun_{K_n}^{(k)}(x)
 \qquad(k\ge2,\ x<0)
\end{equation}
\cite[Conjecture~3]{DongEtAl}.
Ning and Yang proved \eqref{eq:C2} in the range
$x\le-3.01\Delta(G)k$, where $\Delta(G)$ is the maximum degree
\cite[Theorem~1.5]{NingYang}.

For positive integers $s_1,\ldots,s_p$, let
$\ThetaG(s_1,\ldots,s_p)$ be the graph obtained from two terminals by joining
them with $p$ internally vertex-disjoint paths of lengths
$s_1,\ldots,s_p$, measured in edges.

\begin{theorem}\label{thm:main}
Let $H=\ThetaG(8,8,9,9)$. Then $H$ is a simple planar graph on $32$ vertices
and
\[
 \Lfun_H^{(3)}\!\left(-\tfrac12\right)
 =
 \frac{5255579329573016681520193060821596608}
 {18016243181285229121994454910138997625}>0.
\]
Consequently, Conjecture~2 of Dong et al. is false, even for simple planar
graphs.
\end{theorem}

The witness does not conflict with the positive result quoted above. Here
$\Delta(H)=4$ and $k=3$, so \cite[Theorem~1.5]{NingYang} asserts
\eqref{eq:C2} for $H$ and $k=3$ only in the range
$x\le-3.01\cdot4\cdot3=-36.12$, and the witness point $x=-1/2$ lies far
outside it. Theorem~\ref{thm:main} therefore locates the failure of
\eqref{eq:C2} strictly between $-36.12$ and $0$, where nothing was proved.

\medskip
\noindent\textbf{Status.}
What is refuted is \eqref{eq:C2} as a universally quantified statement over
all graphs, all $k\ge2$ and all $x<0$; one exact witness settles that. What is
not touched: the case $k=2$ of \eqref{eq:C2}, for which the present witness
gives a negative value (see the Remark below); the range
$x\le-3.01\Delta(G)k$, which is a theorem; and the lower inequality of
\eqref{eq:C3}, since every spanning tree $Q$ of $H$ has
$\Lfun_Q^{(3)}(-1/2)=-928/27<\Lfun_H^{(3)}(-1/2)$. Only the upper inequality
of \eqref{eq:C3} is refuted here, and only at $k=3$, $x=-1/2$.

\medskip
\noindent\textbf{Statement numbering.}
The labels ``Conjecture~2'' and ``Conjecture~3'' are those of the accepted
version arXiv:1803.08658v3 of \cite{DongEtAl}, which is the only full text we
could reach; the printed journal text was not accessible to us and its
numbering may differ. Both statements are reproduced in full in
\eqref{eq:C2} and \eqref{eq:C3} above, so what is refuted here does not depend
on the labelling.

\section{Exact calculation}

For a path of length $s$ with prescribed endpoint colors, set
\[
 A_s(q)=\frac{(q-1)^s+(q-1)(-1)^s}{q},
 \qquad
 D_s(q)=\frac{(q-1)^s-(-1)^s}{q}.
\]
The numerators vanish at $q=0$, so $A_s,D_s\in\mathbb Z[q]$.

The next formula is not new. After clearing the factors of $q$ in the
$A_{s_i}$ and $D_{s_i}$ it is algebraically identical to
\cite[eq.~(2.6a)]{BHSW}, and for $p=3$ the authors of \cite{BHSW} attribute it
to Read and Tutte. We record a proof only to fix the normalisation used in
Theorem~\ref{thm:main}.

\begin{lemma}[Brown, Hickman, Sokal, and Wagner \cite{BHSW}]\label{lem:theta}
For positive integers $s_1,\ldots,s_p$,
\begin{equation}\label{eq:theta}
 P\bigl(\ThetaG(s_1,\ldots,s_p),q\bigr)
 =q\left(
   \prod_{i=1}^p A_{s_i}(q)
   +(q-1)\prod_{i=1}^p D_{s_i}(q)
 \right).
\end{equation}
\end{lemma}

\begin{proof}
Fix an integer $q\ge2$, let $J$ be the all-ones $q\times q$ matrix and $I$ the
identity. Then $J-I$ is the adjacency matrix of $K_q$, so the $(c,c')$ entry of
$(J-I)^s$ is the number of proper colorings of the internal vertices of a path
of length $s$ whose endpoints are colored $c$ and $c'$. Since $J^2=qJ$,
induction on $s$ gives
\[
 (J-I)^s=D_s(q)\,J+(-1)^sI
\]
(equivalently: $J-I$ has eigenvalues $q-1$ once and $-1$ with multiplicity
$q-1$). Its diagonal entries are $D_s(q)+(-1)^s=A_s(q)$ and its off-diagonal
entries are $D_s(q)$, so that count is $A_s(q)$ when $c=c'$ and $D_s(q)$ when
$c\ne c'$. Now color one terminal of $\ThetaG(s_1,\ldots,s_p)$, in $q$ ways.
The other terminal either receives the same color, the $p$ paths then
contributing $\prod_iA_{s_i}(q)$, or one of the other $q-1$ colors,
contributing $(q-1)\prod_iD_{s_i}(q)$. This gives \eqref{eq:theta} for every
integer $q\ge2$, and since both sides are polynomials in $q$, the identity
follows.
\end{proof}

\begin{proof}[Proof of Theorem~\ref{thm:main}]
The vertex count is
\[
 2+(8-1)+(8-1)+(9-1)+(9-1)=32.
\]
All path lengths are at least two, so $H$ is simple, and the generalized theta
construction is planar. Since $|V(H)|$ is even,
$\Lfun_H(x)=\log P(H,x)$ for $x<0$.

Applying \eqref{eq:theta} to $H$ and differentiating at $q=-1/2$ gives
\begin{align*}
 P\left(H,-\tfrac12\right)
 &=\frac{16528743341798025}{4294967296},
 &
 P'\left(H,-\tfrac12\right)
 &=-\frac{38758016767164565}{536870912},
 \\
 P''\left(H,-\tfrac12\right)
 &=\frac{5549131014505347}{4194304},
 &
 P'''\left(H,-\tfrac12\right)
 &=-\frac{1586794036505230809}{67108864}.
\end{align*}
For any nonvanishing function $f$,
\[
 (\log f)'''
 =\frac{f'''}{f}
 -3\frac{f'f''}{f^2}
 +2\left(\frac{f'}{f}\right)^3.
\]
Substitution yields
\[
 \Lfun_H^{(3)}\!\left(-\tfrac12\right)
 =
 \frac{5255579329573016681520193060821596608}
 {18016243181285229121994454910138997625}>0,
\]
which contradicts \eqref{eq:C2}.
\end{proof}

\begin{remark}
The failure is narrow, and a reader sampling other values will see negative
numbers. At the same point the second derivative is still negative,
\[
 \Lfun_H^{(2)}\!\left(-\tfrac12\right)
 =-\frac{55802572864408523561071168}{6872416789892908604843025}
 =-8.1197\ldots,
\]
so the sign change is a third-order phenomenon; and on the grid of step
$10^{-3}$ in $[-3/5,-3/10]$ the inequality $\Lfun_H^{(3)}(x)>0$ holds exactly
at the points with $-0.518\le x\le-0.433$, the sign being negative at both
neighboring grid points. In particular $\Lfun_H^{(3)}(-1)<0$ and
$\Lfun_H^{(3)}(-1/4)<0$: the witness is specific to $k=3$ and to $x$ near
$-1/2$.
\end{remark}

\begin{corollary}\label{cor:C3}
Conjecture~3 of Dong et al.\ \cite[Conjecture~3]{DongEtAl} is false. More
precisely, its upper inequality fails for $G=H$, $k=3$, and $x=-1/2$.
\end{corollary}

\begin{proof}
For $x<0$,
\[
 (-1)^nP(K_n,x)=\prod_{j=0}^{n-1}(j-x),
 \qquad
 \Lfun_{K_n}^{(k)}(x)
 =-(k-1)!\sum_{j=0}^{n-1}(j-x)^{-k}<0.
\]
Theorem~\ref{thm:main} gives
$\Lfun_H^{(3)}(-1/2)>0$, so the upper inequality in \eqref{eq:C3} fails.
The graph $H$ is connected and noncomplete, and any spanning tree of $H$ is
a chordal proper spanning subgraph, so all hypotheses of Conjecture~3 are
satisfied.
\end{proof}

\section{Reproduction and verification}

\noindent\textbf{Reproduction.}
Lemma~\ref{lem:theta} determines $P(H,q)$ from the parameters $(8,8,9,9)$
alone, so the calculation can be redone in a few lines of exact rational
arithmetic: expand \eqref{eq:theta} with $p=4$ and $(s_i)=(8,8,9,9)$, take the
first three derivatives at $q=-1/2$, and substitute into the identity for
$(\log f)'''$ above. For a concrete labelling, take the terminals to be $0$
and $1$ and number the internal vertices consecutively along the four paths,
\begin{center}
\texttt{0-2-3-4-5-6-7-8-1},\qquad
\texttt{0-9-10-11-12-13-14-15-1},\\
\texttt{0-16-17-18-19-20-21-22-23-1},\qquad
\texttt{0-24-25-26-27-28-29-30-31-1},
\end{center}
which are the $34$ edges of $H$; the corresponding graph6 string is the
concatenation of the two lines
\begin{center}
\texttt{\_PCGGD@\_??\_@?@??\_?I?@\_????G??G??C??@???G???g??@\_???????\_???G???@}\\
\texttt{????C????G????I????C}
\end{center}

\medskip
\noindent\textbf{Exact verification.}
The computation reported here was re-run independently, on a separate machine
and from the graph6 string printed above, in exact integer and rational
arithmetic throughout. It decodes and re-encodes
$H$; confirms by an explicit bijection that the decoded graph is
$\ThetaG(8,8,9,9)$ with $32$ vertices, $34$ edges and $\Delta(H)=4$; checks
planarity twice, by a Kuratowski branch-vertex count and by tracing the faces
of an explicit rotation system against Euler's formula; computes $P(H,q)$ both
from \eqref{eq:theta} and, independently, by a generic proper-coloring count at
$q=0,\ldots,32$ followed by exact interpolation; recomputes the four
derivatives displayed above and $\Lfun_H^{(3)}(-1/2)$ by two routes; and
rechecks the hypotheses of \eqref{eq:C3}, including that all $2448$ spanning
trees of $H$ are chordal proper spanning subgraphs. All $80$ of its checks
passed, and every value it recomputed agreed with the value printed here. The
program used for that re-run, together with its transcript, accompanies this
note; it is not needed in order to redo the calculation, since every input it
consumed is printed above.

One limit of that re-run should be recorded. It also ran a census of
generalized theta graphs, which is not among the claims of this note and which
the re-run itself records as bounded: the census is complete only for order at
most $34$, larger orders were not enumerated, and no graph outside the generalized
theta family was examined. Nothing here, therefore, asserts that $H$ is a
counterexample of least order, either among generalized theta graphs or among
graphs.

\begin{thebibliography}{9}

\bibitem{BHSW}
J.~I. Brown, C.~A. Hickman, A.~D. Sokal, and D.~G. Wagner,
\emph{On the chromatic roots of generalized theta graphs},
J. Combin. Theory Ser. B \textbf{83} (2001), no.~2, 272--297.
\href{https://doi.org/10.1006/jctb.2001.2057}{doi:10.1006/jctb.2001.2057};
\href{https://arxiv.org/abs/math/0012033}{arXiv:math/0012033}.

\bibitem{DongEtAl}
F.~M. Dong, J.~Ge, H.~L. Gong, B.~Ning, Z.~Ouyang, and E.~G. Tay,
\emph{Proving a conjecture on chromatic polynomials by counting the number of
acyclic orientations},
J. Graph Theory \textbf{96} (2021), no.~3, 343--360.
\href{https://doi.org/10.1002/jgt.22617}{doi:10.1002/jgt.22617}.
Statement numbers are quoted from the accepted version
\href{https://arxiv.org/abs/1803.08658v3}{arXiv:1803.08658v3}, the only full
text we could access.

\bibitem{NingYang}
B.~Ning and Y.~Yang,
\emph{On derivatives and higher-order derivatives of chromatic polynomials},
arXiv:2604.13221v2, 2026.
\href{https://arxiv.org/abs/2604.13221}{arXiv:2604.13221}.

\end{thebibliography}

\end{document}